% Compile twice with pdfLaTeX from this directory.
% English slides, with speaker notes and timing in each frame.
% Define \presenternotes before input to produce a notes-only PDF.
\documentclass[aspectratio=1610,xcolor=dvipsnames,t]{beamer}
\usepackage{lmodern}
\usepackage{sansmath}
\usepackage{mathastext}
\usepackage{graphicx}
\usepackage{tikz}
\usetikzlibrary{decorations.pathreplacing}
\usepackage{beamerthemesplit}
% Reuse the theme stored in the sibling template directory.
\makeatletter
\def\input@path{{../模板/}}
\makeatother
\usetheme[compress]{Heidelberg}
\definecolor{unirot}{rgb}{0.5976525,0,0}
\usecolortheme[named=unirot]{structure}
\usepackage{amsmath,amssymb}
\newcommand{\toprule}{\hline}
\newcommand{\midrule}{\hline}
\newcommand{\bottomrule}{\hline}
\setbeamertemplate{headline}{}
\ifdefined\presenternotes
  \setbeameroption{show only notes}
\else
  \setbeameroption{hide notes}
\fi
\newcommand{\keyterm}[1]{\textcolor{unirot}{\textbf{#1}}}
% Keep existing backup content while suppressing per-slide source labels.
\newcommand{\papersource}[1]{}
\newcommand{\paperfig}[4][.50\textheight]{%
  \includegraphics[page=#2,viewport=#3,clip,width=#4,height=#1,keepaspectratio]{../2408.03314v1.pdf}}
\title[Test-Time Compute in LLMs]{Scaling LLM Test-Time Compute Optimally can be More Effective than Scaling Model Parameters}
\subtitle{Allocating inference compute under a fixed budget}
\author[Zhenbang Lyu]{Zhenbang Lyu}
\institute[Uni HD]{Heidelberg University\\Institute of Computer Science}
\date[]{Charlie Snell, Jaehoon Lee, Kelvin Xu, and Aviral Kumar\\2024, arXiv:2408.03314v1}

\begin{document}

\begin{frame}[plain]
  \titlepage
  \note{Time 0:00--0:30. The paper asks how to improve a fixed model by spending more computation when it answers a question. We will first examine how to allocate that computation, then compare it with training a larger model.}
\end{frame}

\begin{frame}{Can Language Models Benefit from More Thinking?}
  Humans often think for longer about difficult problems to reach better answers and improve their decisions.
  \par\medskip
  Extra thought gives us time to reconsider an initial answer and work through the reasoning more carefully.
  \par\medskip
  \keyterm{Can language models also improve their answers by using more computation at test time?}
  \par\medskip
  The paper studies how a trained model can use an additional inference budget on a challenging question.
  \note{Time 0:30--1:15. Follow the introduction's motivation. Humans tend to spend longer thinking about difficult problems to improve their decisions. Ask whether a trained language model can gain a similar ability. For a model, the resource we can allocate is additional computation during inference. The paper asks how effectively this resource can improve accuracy on a given prompt. Next make the available uses of that computation concrete.}
\end{frame}

\begin{frame}{How Do We Work through a Math Problem?}
  Suppose we want a better chance of finding the correct answer to a math problem.
  \par\medskip
  \begin{itemize}
    \item \textbf{Ask several people to solve it independently.} Collect their complete solutions, then compare the answers.
    \par\smallskip
    \item \textbf{Build on an existing attempt.} Read someone else's working and use it to develop an improved solution.
    \par\smallskip
    \item \textbf{Ask for feedback halfway through.} Show someone our unfinished reasoning before deciding how to continue.
  \end{itemize}
  \par\medskip
  \keyterm{A model can also try independently, revise earlier work, or use feedback during generation.}
  \note{Time 1:15--2:00. Start with several people independently solving the same problem. This motivates parallel sampling and best-of-N as the simplest baseline. Then introduce two ways to use existing work: revise an earlier attempt or obtain feedback before finishing. In the paper, parallel samples and revision history come from the model itself.}
\end{frame}

\begin{frame}{Two Directions for Using Extra Computation}
  \textbf{Baseline}\quad Generate $N$ independent solutions, then select an answer with a verifier: \keyterm{best-of-$N$}.
  \par\medskip
  The paper improves on this baseline in two directions.
  \par\smallskip
  \begin{itemize}
    \item \textbf{Change what the model generates.} Give it earlier answers and ask for a revised solution. \keyterm{Revisions} change its \textbf{proposal distribution}.
    \par\smallskip
    \item \textbf{Change how feedback guides generation.} Score partial solutions and choose which paths to develop. This is \keyterm{verifier-guided search}.
  \end{itemize}
  \par\medskip
  We first follow the verifier-guided search route from intermediate steps to a final answer.
  \note{Time 2:00--3:00. Map independent solvers to independent model samples. Best-of-N generates complete candidates and uses a verifier for final selection. The two improvements correspond to the proposer/verifier framework. Revisions condition on previous attempts. Search uses partial-solution scores. All share a finite budget. Focus on the verifier-guided route first and show the complete workflow before explaining its components.}
\end{frame}

\begin{frame}{PRM-Guided Search from Steps to an Answer}
  Score partial solutions and stop low-scoring branches before they become complete answers.
  \par\smallskip
  \begin{center}
  \resizebox{.96\textwidth}{!}{%
    \begin{tikzpicture}[x=1cm,y=.8cm,>=stealth,
      prefixnode/.style={draw=unirot,rounded corners=2pt,fill=unirot!4,minimum width=1.15cm,minimum height=.43cm,font=\small},
      stopped/.style={prefixnode,draw=black!25,fill=black!4,text=black!40},
      answernode/.style={draw=unirot,minimum width=.7cm,minimum height=.43cm,font=\small},
      flow/.style={->,line width=.7pt,draw=black!70},
      cutflow/.style={->,line width=.7pt,draw=black!30}]
      \node[font=\small] (q) at (0,-.25) {Question};
      \foreach \i/\h in {1/1.2,2/0,3/-1.2}{
        \node[prefixnode] (a\i) at (2.1,\h) {Step 1};
        \draw[flow] (q) -- (a\i);
      }
      \node[stopped] (a4) at (2.1,-2.3) {Step 1};
      \draw[cutflow] (q) -- (a4);
      \node[text=black!50,font=\large] at (2.95,-2.3) {$\times$};
      \foreach \i/\keep/\drop in {1/1.5/.85,2/.25/-.4,3/-1/-1.65}{
        \node[prefixnode] (b\i) at (4.9,\keep) {Step 2};
        \node[stopped] (c\i) at (4.9,\drop) {Step 2};
        \draw[flow] (a\i) -- (b\i);
        \draw[cutflow] (a\i) -- (c\i);
        \node[text=black!50,font=\large] at (5.75,\drop) {$\times$};
        \node (dots\i) at (6.7,\keep) {$\cdots$};
        \node[answernode] (y\i) at (8.1,\keep) {$Y_{\i}$};
        \draw[flow] (b\i) -- (dots\i);
        \draw[flow] (dots\i) -- (y\i);
      }
      \node[font=\small,text=unirot] at (4.4,2.2) {PRM scoring and pruning after each expansion};
      \node[draw=unirot,rounded corners=2pt,align=center,font=\small,minimum height=.72cm] (select) at (10.0,.25) {Answer\\selection};
      \node[font=\large,text=unirot] (final) at (12.0,.25) {$\hat{Y}$};
      \foreach \i in {1,2,3}{\draw[flow] (y\i) -- (select);}
      \draw[flow] (select) -- (final);
      \node[font=\small,text=black!60] at (7.6,-2.3) {$\times$ Pruned: no further generation};
    \end{tikzpicture}%
  }
  \end{center}
  \small
  Repeat until completion or the expansion limit. Three decisions remain:
  \begin{enumerate}
    \item Who scores the reasoning?
    \item How do we select the final answer?
    \item Which candidates should the search retain and expand?
  \end{enumerate}
  \note{Time 3:00--4:00. Follow the branching tree from the question. Gray nodes marked with a cross are pruned using PRM scores and receive no further generation. Retained Step 1 prefixes branch into different Step 2 candidates, which are scored and pruned again. Only surviving paths reach complete answers Y1, Y2, and Y3 for final selection. Standard best-of-N completes all independent candidates before selecting. Here partial-solution scores decide which branches receive later compute. The diagram is schematic and does not specify a fixed beam size. The next slides explain the scorer, final selection, and search rule.}
\end{frame}

\begin{frame}{Who Can Score the Reasoning?}
  Train a \keyterm{process reward model (PRM)} to predict the chance of finishing correctly from a given prefix.
  \par\smallskip
  \begin{center}
  \resizebox{.98\textwidth}{!}{%
    \begin{tikzpicture}[x=1cm,y=1cm,>=stealth,
      future/.style={draw=black!55,dashed,rounded corners=2pt,minimum height=.55cm,font=\small,align=center},
      rollout/.style={->,dashed,draw=black!55,line width=.7pt}]
      \node[draw=unirot,line width=.9pt,rounded corners=2pt,fill=unirot!4,
        align=center,minimum width=2.45cm,minimum height=1.1cm,font=\small] (prefix) at (1.3,0)
        {Question\\+ Steps $1,\ldots,t$};
      \node[font=\small,text=unirot] at (1.3,.95) {Current prefix};
      \node[font=\small] at (6.0,2.45) {Sample four continuations to complete answers};
      \foreach \i/\h in {1/1.8,2/.6,3/-.6,4/-1.8}{
        \node[future,minimum width=1.6cm] (next\i) at (4.6,\h) {Next steps};
        \node[future,minimum width=1.9cm] (answer\i) at (7.5,\h) {Final answer};
        \draw[rollout] (prefix.east) -- (next\i.west);
        \draw[rollout] (next\i.east) -- node[midway,fill=white,inner sep=2pt] {$\cdots$} (answer\i.west);
      }
      \foreach \i/\h in {1/1.8,2/.6,3/-.6}{
        \node[font=\Large,text=ForestGreen] at (8.85,\h) {$\checkmark$};
      }
      \node[font=\Large,text=unirot] at (8.85,-1.8) {$\times$};
      \draw[decorate,decoration={brace,amplitude=7pt},line width=.8pt,draw=black!65]
        (9.35,2.12) -- (9.35,-2.12);
      \node[align=center,font=\small] at (11.15,.95) {3 correct out of 4};
      \node[font=\Large,text=unirot] at (11.15,.1) {$\frac{3}{4}=0.75$};
      \node[align=center,font=\small,text=unirot] at (11.15,-.85) {Training target for\\the current prefix};
      \node[font=\small,text=black!60] at (6.3,-2.65) {Illustrative training rollouts. Dashed paths are possible futures.};
    \end{tikzpicture}%
  }
  \end{center}
  \small
  During search, the trained PRM \textbf{scores new prefixes directly}.
  \note{Time 4:00--5:15. Point to the solid box, which contains the question and the entire prefix through the current step. Hold that prefix fixed and independently sample four continuations with the base model, each all the way to a final answer. All future nodes and connecting paths are dashed. Check final correctness using ground truth during training. Three checks and one cross give an illustrative target of 3/4=0.75 for the current prefix. This estimates eventual success under the base model's sampling policy, not whether an isolated step is logically correct. Train the PRM to predict such targets from many prefixes. At search time the learned PRM directly scores a prefix without repeating these training rollouts. Four is an illustrative sample count, not the experiment's rollout count.}
\end{frame}

\begin{frame}{How Do We Select the Final Answer?}
  Use the PRM's \textbf{last-step score} for each complete solution, then sum scores for matching final answers.
  \par\medskip
  \begin{center}
  \begin{tabular}{lcc}
    \toprule
    Example candidate & Final answer & Last-step score\\
    \midrule
    $Y_1$ & $x=2$ & 0.60\\
    $Y_2$ & $x=2$ & 0.55\\
    $Y_3$ & $x=3$ & 0.90\\
    \bottomrule
  \end{tabular}
  \end{center}
  \[
    \operatorname{score}(x=2)=1.15,\qquad
    \operatorname{score}(x=3)=0.90
  \]
  \keyterm{Weighted best-of-$N$ selects $x=2$.}
  \par\smallskip
  Standard best-of-$N$ would pick the single highest-scoring candidate. Majority voting counts answers without weights.
  \note{Time 5:15--6:15. Answer the second workflow question using the three endpoints in the overview. The final-step PRM score sees the whole completed prefix. One aggregation turns step scores into a solution score; the other pools candidate scores by final answer. Candidate Y3 has the highest individual score, but x=2 has more total weighted support. A sum of 1.15 is not a probability. All numbers are illustrative. This is the final selection stage shared by the search methods. Next ask how the candidates themselves should be generated.}
\end{frame}

\begin{frame}{Which Candidates Should We Develop?}
  \begin{center}
    \paperfig[.53\textheight]{8}{61 493 551 710}{.98\textwidth}
  \end{center}
  \small
  \textbf{Best-of-$N$}\quad Generate independent complete solutions, then select.\\
  \textbf{Beam search}\quad Score prefixes, keep the top candidates, then expand.\\
  \textbf{Lookahead}\quad Temporarily continue $k$ steps before scoring a prefix.
  \par\smallskip
  Lookahead spends additional compute on its temporary continuations.
  \note{Time 6:15--7:45. Answer the third workflow question. Best-of-N waits until candidates are complete. Beam search allocates generation toward promising prefixes. Lookahead uses a temporary greedy continuation of up to k steps and scores its endpoint to evaluate the current prefix. The temporary steps help rank current prefixes. k=0 gives ordinary beam search. The paper estimates lookahead cost as N times k+1, so fair comparisons include that overhead. The detailed numerical example is in backup. Next compare these methods at matched compute budgets.}
\end{frame}


\begin{frame}{Beam Search Helps Most at Low Budgets}
  {\small PaLM 2-S* on MATH: 12,000 training questions and 500 test questions.}
  \par\medskip
  \begin{columns}[T,onlytextwidth]
    \begin{column}{.58\textwidth}
      \centering
      \paperfig[.61\textheight]{9}{61 539 309 708}{\linewidth}
    \end{column}
    \begin{column}{.39\textwidth}
      \small
      \textbf{Beam search is strongest at low budgets, but its advantage shrinks as compute increases.}
      \par\medskip
      At matched cost, extra lookahead usually gives less benefit.
    \end{column}
  \end{columns}
  \note{Time 7:45--9:00. Read Figure 3 left by comparing methods vertically at the same budget. Follow the beam-search curves from the low-budget region to higher budgets and compare with weighted best-of-N. Explain that lookahead's extra rollouts are charged to the budget. State both displayed conclusions. Then ask whether the aggregate trend hides different outcomes across question difficulty.}
\end{frame}

\begin{frame}{Difficulty Changes Which Search Strategy Wins}
  Each difficulty bin has four budgets, left to right: \textbf{4, 16, 64, 256}.
  \par\medskip
  \begin{columns}[T,onlytextwidth]
    \begin{column}{.52\textwidth}
      \centering
      \paperfig[.55\textheight]{9}{315 539 551 708}{\linewidth}
      \par\smallskip
      \small Read the top of each color. The bars overlap.
      \par\medskip
      \raggedright
      \small\textbf{Example, bin 2:} beam wins at budget 4. Best-of-$N$ wins at 256.
    \end{column}
    \begin{column}{.45\textwidth}
      \small
      \keyterm{Easy questions (bins 1--2)}
      \par\smallskip
      Beam search can win at small budgets, then stall or decline. Best-of-$N$ catches up.
      \par\medskip
      \keyterm{Medium difficulty (bins 3--4)}
      \par\smallskip
      Beam search consistently beats best-of-$N$ at the same budget.
      \par\medskip
      \keyterm{Hardest questions (bin 5)}
      \par\smallskip
      All three methods make little progress.
      \par\medskip
      \keyterm{Choose the search setting using both difficulty and budget.}
    \end{column}
  \end{columns}
  \note{Time 9:00--11:15. Figure 3 right and the paragraph Which problems does search improve? in Section 5.3 explicitly discuss these findings. Read the four positions within each difficulty bin as budgets 4, 16, 64, and 256. Blue is beam search, orange is weighted best-of-N, and green is majority voting. All bars start at zero and overlap, so compare the top of each color rather than adding the visible segments. In bin 2, compare blue and orange at budget 4, then at budget 256: the winner changes from beam search to best-of-N. In easy bins 1--2, beam search can improve initially and later stall or decline as the budget grows; the authors identify signs of over-optimization. This does not mean beam search always loses on easy questions or decreases at every budget increment. In bins 3--4, beam search consistently beats best-of-N at matched budgets. In bin 5, none makes meaningful progress. The authors suggest that stronger optimization may exploit imperfect PRM scores on easy questions, while on medium questions search helps guide a base model that rarely samples correct solutions. Treat this as their proposed explanation. The displayed conclusion motivates selecting the search setting by difficulty and budget in the next slide. Formal bin estimation remains in backup.}
\end{frame}

\begin{frame}{Difficulty-Based Allocation Improves Search Efficiency}
  \begin{columns}[T,onlytextwidth]
    \begin{column}{.58\textwidth}
      \centering
      \paperfig[.63\textheight]{10}{306 466 551 635}{\linewidth}
    \end{column}
    \begin{column}{.39\textwidth}
      \small
      \textbf{Difficulty-based allocation improves search efficiency.}
      \par\medskip
      At low budgets, 16 generations nearly match best-of-$N$ with 64.
      \par\medskip
      Predicted difficulty loses some of its advantage at higher budgets.
    \end{column}
  \end{columns}
  \par\smallskip
  {\small These generation budgets exclude difficulty-estimation cost.}
  \note{Time 11:15--12:00. Compare the adaptive blue/orange curves with PRM weighted best-of-N in green. Compare at the same accuracy across budgets to illustrate 16 versus 64. Distinguish true-correctness oracle bins from verifier-predicted bins. The result is a low-budget saving, and estimating difficulty is excluded. Then return to the second route: revising full answers.}
\end{frame}

\begin{frame}{Revising Answers and Selecting the Output}
  A fine-tuned revision model uses earlier full solutions to write a new answer.
  \par\medskip
  \textbf{1\quad Generate}\quad Each call sees the question and up to four most recent answers.
  \begin{center}
  \resizebox{.97\textwidth}{!}{%
    \begin{tikzpicture}[x=1cm,y=1cm,>=stealth,
      revanswer/.style={minimum width=.86cm,minimum height=.58cm,font=\large},
      revflow/.style={->,line width=.75pt,draw=black!65}]
      \filldraw[fill=unirot!6,draw=unirot!65,rounded corners=3pt]
        (3.93,-.36) rectangle (8.97,.36);
      \node[revanswer,text=black!55] (a1) at (.6,0) {$A_1$};
      \node[revanswer,text=black!55] (a2) at (1.9,0) {$A_2$};
      \node[revanswer,text=black!55] (dots) at (3.2,0) {$\cdots$};
      \node[revanswer] (a60) at (4.5,0) {$A_{60}$};
      \node[revanswer] (a61) at (5.8,0) {$A_{61}$};
      \node[revanswer] (a62) at (7.1,0) {$A_{62}$};
      \node[revanswer] (a63) at (8.4,0) {$A_{63}$};
      \node[revanswer,draw=unirot,rounded corners=3pt,text=unirot] (a64) at (10.7,0) {$A_{64}$};
      \draw[revflow] (a1) -- (a2);
      \draw[revflow] (a2) -- (dots);
      \draw[revflow] (dots) -- (a60);
      \draw[revflow] (a60) -- (a61);
      \draw[revflow] (a61) -- (a62);
      \draw[revflow] (a62) -- (a63);
      \draw[revflow] (a63) -- (a64);
      \node[font=\small,text=black!55] at (1.9,.88) {Earlier answers};
      \node[font=\small,text=unirot] at (6.45,.88)
        {Input for $A_{64}$: $q+A_{60}+A_{61}+A_{62}+A_{63}$};
      \node[font=\small,text=unirot] at (6.45,-.75) {The four-answer context window};
      \node[font=\small,text=black!60] at (10.7,-.75) {New answer};
    \end{tikzpicture}%
  }
  \end{center}
  \textbf{2\quad Select}\quad Keep all candidates, then use a verifier to select an answer.
  \begin{center}
  \resizebox{.96\textwidth}{!}{%
    \begin{tikzpicture}[x=1cm,y=1cm,>=stealth,
      pickflow/.style={->,line width=.75pt,draw=black!65}]
      \node[font=\small,text=black!60] at (1.9,1.0) {All candidates};
      \node[font=\small,text=black!60] at (5.8,1.0) {One answer per chain};
      \node[font=\large] (pa) at (1.9,.4) {$\{A_1,\ldots,A_{64}\}$};
      \node[font=\large] (pb) at (1.9,-.4) {$\{B_1,\ldots,B_{64}\}$};
      \node[font=\large] (ya) at (5.8,.4) {$\hat A$};
      \node[font=\large] (yb) at (5.8,-.4) {$\hat B$};
      \node[draw=unirot,rounded corners=3pt,font=\small,align=center,
        minimum width=1.8cm,minimum height=.85cm] (across) at (8.5,0) {Select across\\chains};
      \node[font=\Large,text=unirot] (y) at (10.7,0) {$\hat y$};
      \draw[pickflow] (pa) -- (ya);
      \draw[pickflow] (pb) -- (yb);
      \draw[pickflow] (ya) -- (across);
      \draw[pickflow] (yb) -- (across);
      \draw[pickflow] (across) -- (y);
    \end{tikzpicture}%
  }
  \end{center}
  {\small Training uses incorrect references. At inference, earlier answers may already be correct, and a later revision can undo them.}
  \note{Time 12:00--13:15. First introduce a model fine-tuned to use previous full answers. The first call sees only the question. The next sees the question plus A1. At round 64 it sees the question and answers 60 through 63, highlighted in red. The older answers remain stored even though they have left the input window. Next follow the lower diagram: every answer in a chain is available for verifier-based selection, then selected answers can be compared across chains. Two chains of length 64 imply 128 candidates in total; the diagram is illustrative and other chain lengths are possible. With one chain, the paper also evaluates majority voting. The selected output need not be the last answer because revisions can undo correct answers. Training-pair construction and the reported conversion rate are in backup. This sets up Figure 6, which separately measures quality at each revision step and accuracy after final selection.}
\end{frame}

\begin{frame}{Revision Quality and Final Answer Accuracy}
  \begin{center}
    \paperfig[.43\textheight]{12}{61 539 551 715}{.98\textwidth}
  \end{center}
  \small
  \begin{columns}[T,onlytextwidth]
    \begin{column}{.48\textwidth}
      \textbf{Left figure: the answer at round $t$}
      \par\smallskip
      \textbf{One question: 4 chains, 64 answers per chain.}
      \par\smallskip
      Average round-$t$ correctness over the four chains, then over all \textbf{500 test questions}.
      \par\smallskip
      \keyterm{Most improvement happens in early rounds.}
    \end{column}
    \begin{column}{.48\textwidth}
      \textbf{Right figure: one selected answer}
      \par\smallskip
      Per question: $N$ independent answers or $N$ sequential answers.
      \par\smallskip
      Select one from $N$, then measure accuracy across all \textbf{500 test questions}.
      \par\smallskip
      \keyterm{Sequential is slightly better with either selector.}
    \end{column}
  \end{columns}
  \note{Time 13:15--15:00. Both figures summarize all 500 test questions. For the left figure, first explain the data for one question: four independent revision chains, each containing 64 complete answers. At round 10, check the tenth answer in each of the four chains. Repeat for all questions and average. This checks the answer at a fixed round without selecting the best earlier answer. Most improvement happens early and later accuracy changes only a little. For the right figure, use N=8: the same revision model either generates eight independent answers or an eight-answer sequential chain for each question. A selector chooses one from those eight, and the chosen answer is checked. Average over all questions. Two generation styles times two selectors make four curves. Compare sequential weighted selection with parallel weighted selection, then sequential majority with parallel majority. Sequential is slightly better in both comparisons. The four chains of length 64 describe the left figure's measurement, not four curves or four questions in the right figure.}
\end{frame}

\begin{frame}{The Best Mix Depends on Difficulty and Budget}
  \begin{center}
    \paperfig[.44\textheight]{13}{61 539 551 708}{.98\textwidth}
  \end{center}
  \small
  \begin{itemize}
    \item At high budgets, a mixture can beat either pure parallel or pure sequential generation.
    \item Easy questions favor long chains; harder questions often benefit from more independent starts.
  \end{itemize}
  \par\smallskip
  \keyterm{The best mix depends on both difficulty and budget.}
  \note{Time 15:00--16:45. The left x-axis is chain length divided by chain count. A line keeps their product, the total budget, fixed. Color is budget on the left and ratio on the right. The right figure fixes 128 generations and groups by difficulty. Read a peak within a fixed-budget curve and then compare colors within difficulty bins. State all three displayed conclusions.}
\end{frame}

\begin{frame}{Adaptive Revisions Use the Generation Budget More Efficiently}
  \begin{columns}[T,onlytextwidth]
    \begin{column}{.58\textwidth}
      \centering
      \paperfig[.63\textheight]{13}{306 186 551 357}{\linewidth}
    \end{column}
    \begin{column}{.39\textwidth}
      \small
      \textbf{Adaptive revisions keep improving after parallel sampling begins to plateau.}
      \par\medskip
      \keyterm{64 candidates can beat parallel best-of-$N$ with 256 candidates.}
      \par\medskip
      Both oracle and predicted difficulty improve the allocation.
    \end{column}
  \end{columns}
  \par\smallskip
  {\small The reported saving excludes difficulty-estimation cost.}
  \note{Time 16:45--17:45. Figure 8 compares adaptive mixtures with pure parallel sampling. Solid lines use weighted best-of-N and dashed lines use majority. Colors distinguish oracle, predicted, and parallel. Compare within a selection rule. Read across budgets for the 64-versus-256 example. State the displayed conclusions and the excluded difficulty cost. Next compare extra inference with larger model size at equal total FLOPs.}
\end{frame}

\begin{frame}{Pretraining and Inference at Equal Total FLOPs}
  Extra inference improves the small model. How does it compare with training a larger model?
  \par\smallskip
  Hold training tokens fixed and increase parameters by a factor $M$.
  \[
    X\approx6PD_{\mathrm{pretrain}},\qquad
    Y\approx2PD_{\mathrm{inference}}
  \]
  $P$ is the parameter count. $D_{\mathrm{inference}}$ is the total inference token workload.
  \par\medskip
  \[
    \underbrace{X+aY}_{\text{small model, extra inference}}
    =\underbrace{M(X+Y)}_{\text{large model, greedy decoding}}
  \]
  \par\medskip
  The larger model costs more both to train and to run. The small model can spend its training savings on inference.
  \note{Time 17:45--19:15. Move from choosing inference strategies to choosing where to spend the total compute budget. P denotes parameter count, distinct from earlier candidate count N. X and Y are the small model's baseline training and inference costs. The larger model costs M times as much for both, with training tokens fixed. The small model can use the training savings for more inference. We next calculate how much extra inference that buys.}
\end{frame}

\begin{frame}{Expected Use Determines the Extra Inference Budget}
  \small
  $R$ compares all planned inference tokens with the tokens used for pretraining:
  \[
    R=\frac{D_{\mathrm{inference}}}{D_{\mathrm{pretrain}}}.
  \]
  From $X+aY=M(X+Y)$, the small model's inference multiplier is
  \[
    \begin{aligned}
      a &= M+(M-1)\frac{X}{Y}\\
        &= M+(M-1)\frac{6P D_{\mathrm{pretrain}}}{2P D_{\mathrm{inference}}}\\
        &= M+3(M-1)\frac{D_{\mathrm{pretrain}}}{D_{\mathrm{inference}}}
         = M+\frac{3(M-1)}{R}.
    \end{aligned}
  \]
  \begin{columns}[T,onlytextwidth]
    \begin{column}{.48\textwidth}
      \keyterm{Low $R$: relatively little planned use}
      \par\smallskip
      Example: generate a small batch of self-improvement data.
      \par\smallskip
      More extra compute per answer.
    \end{column}
    \begin{column}{.48\textwidth}
      \keyterm{High $R$: heavy planned use}
      \par\smallskip
      Example: a public chatbot serving many users over time.
      \par\smallskip
      Less extra compute per answer.
    \end{column}
  \end{columns}
  \note{Time 19:15--20:30. R compares the total expected inference token workload with pretraining tokens. It covers all planned use, not the length of one answer; define it for the baseline workload before adding extra test-time compute. Keep the previous slide's equal-total-cost condition. Solve for a to get M+(M-1)X/Y. Substitute the cost estimates: six divided by two gives three, parameters cancel, and Dpretrain/Dinference equals 1/R. Thus a=M+3(M-1)/R. Here a is how many times the baseline inference compute the small model can afford at the matched total cost. With relatively little planned use, such as a small batch of model-improvement data, training savings support more compute per answer. With heavy planned use, such as a large chatbot service, the same savings must support many answers. The three numerical workloads are compared later, so no table is repeated here.}
\end{frame}

\begin{frame}{Extra Inference Can Beat a Larger Model in Some Settings}
  \begin{center}
    \paperfig[.47\textheight]{15}{64 526 549 715}{.98\textwidth}
  \end{center}
  \small
  \begin{itemize}
    \item Extra inference often wins on easy or medium questions, especially when expected inference use is low.
    \item A larger model is often better on very hard questions or when expected inference use is high.
  \end{itemize}
  \par\smallskip
  \keyterm{The comparison depends on both difficulty and the expected inference workload.}
  \note{Time 20:30--22:45. Read Figure 9. Curves are adaptive inference with the small model, stars are greedy answers from the model with about 14 times more parameters. Match color and the workload's vertical line. A curve higher than the star favors additional inference. The same star height repeats at different matched budgets. Purple easy questions and blue hardest questions show different patterns. State both displayed conditional findings. The next slide follows one medium-difficulty group through the three R values. Figure 9 uses oracle difficulty bins.}
\end{frame}

\begin{frame}{Expected Use Can Change Which Model Is Better}
  One example: the \textbf{revision model on difficulty bin 3}, read from the previous figure.
  \par\medskip
  \begin{center}
  \small
  \begin{tabular}{ccl}
    \toprule
    $R$ & \shortstack{Small-model inference\\budget multiplier} & Matched-cost outcome\\
    \midrule
    $22$ & $\sim16\times$ & The larger model is more accurate.\\[2mm]
    $0.79$ & $\sim64\times$ & The two models are close.\\[2mm]
    $0.16$ & $\sim256\times$ & The small model with revisions is more accurate.\\
    \bottomrule
  \end{tabular}
  \end{center}
  \par\medskip
  Lower expected use lets the small model spend more of its training savings on each inference.
  \par\medskip
  \keyterm{The better choice depends on both question difficulty and expected use.}
  \note{Time 22:45--23:45. Follow the green revision curve for difficulty bin 3 in Figure 9. At the high-workload R of 22, the matched budget is near 16 and the larger model's star is higher. Around R=0.79 the matched budget is near 64 and the two are close. At R=0.16 the matched budget is near 256 and the revision curve is higher. These are approximate plot locations for the rounded R values. This is an example for one group and method. The broader conditional conclusions remain those on the preceding slide.}
\end{frame}

\begin{frame}{Limits of the Comparison}
  \begin{itemize}
    \item Estimating difficulty needs substantial compute, excluded from reported budgets.
    \item Hardest problems gain little from the tested search and revision methods.
    \item The evidence comes from PaLM 2-S* on MATH, with capability-specific fine-tuning.
    \item The FLOPs comparison scales parameters at fixed training data and uses oracle difficulty bins.
  \end{itemize}
  \note{Time 23:45--24:30. Real deployment must pay for difficulty estimation. The data and model constrain generality. The larger-model comparison is one parameter-scaling setting. The paper leaves the combination of PRM tree search with revision chains for future study.}
\end{frame}

\begin{frame}{Conclusions}
  \begin{itemize}
    \item The best use of inference compute depends on both the question and the budget.
    \item Difficulty-based allocation improves PRM search and sequential/parallel revisions, with up to fourfold generation-budget savings in reported comparisons.
    \item Extra inference can substitute for larger models on some problems and workloads. The hardest problems remain a limitation.
  \end{itemize}
  \par\medskip
  \keyterm{The title's claim is conditional on problem difficulty and expected inference workload.}
  \note{Time 24:30--25:00. Return to the opening question. A fixed generation budget can be spent in very different ways. Choosing adaptively improves efficiency. The larger-model claim follows only under certain problem and workload conditions. End here at 25 minutes, excluding questions.}
\end{frame}

\begin{frame}{Questions}
  \vfill
  \begin{center}
    {\Huge\color{unirot}\textbf{Thank you}}\\[.4cm]
    {\Large Questions and Discussion}
  \end{center}
  \vfill
  \note{Questions begin after the 25-minute talk. Backup slides contain details on validation, rollouts, and FLOPs accounting.}
\end{frame}

\appendix

\begin{frame}{Backup: Constructing Revision Training Data}
  \begin{columns}[T,onlytextwidth]
    \begin{column}{.43\textwidth}
      \textbf{Build a training pair}
      \par\medskip
      \begin{itemize}
        \item Sample 64 independent answers to one question.
        \item Choose a correct answer as the target.
        \item Add 0--4 related incorrect answers as context.
        \item Prefer close pairs using character edit distance.
      \end{itemize}
    \end{column}
    \begin{column}{.53\textwidth}
      \textbf{An illustrative pair}
      \par\medskip
      Question: solve $2x+3=7$.
      \par\medskip
      \textbf{Incorrect reference}
      \[
        2x=4,\qquad x=4
      \]
      \textbf{Correct target}
      \[
        2x=4,\qquad \textcolor{unirot}{x=2}
      \]
      The target keeps the valid setup and fixes the final calculation.
    \end{column}
  \end{columns}
  \par\medskip
  Fine-tune on \textbf{question + incorrect references} as input, with the \textbf{correct solution} as the target.
  \note{Optional backup. The example is illustrative. In the experiment, 64 complete answers per question are sampled independently at a higher temperature and assembled into synthetic revision trajectories afterward. The number of incorrect reference answers is sampled uniformly from zero through four. Character edit distance prioritizes related incorrect and correct solutions. In the example the setup is shared and the correct target repairs the final calculation. This encourages using earlier work. Edit distance is a heuristic for relatedness and is used to build training pairs. At inference, the window is the most recent four answers in chronological order.}
\end{frame}

\begin{frame}{Backup: Why the Last Revision Can Be Wrong}
  \begin{tabular}{@{}ll@{}}
    \textbf{Training history} & Incorrect references, followed by a correct target.\\[1mm]
    \textbf{Inference history} & Generated references can already be correct.
  \end{tabular}
  \par\medskip
  A possible revision chain, illustrated on $2x+3=7$:
  \begin{center}
    \begin{tikzpicture}[x=1cm,y=1cm,>=stealth,
      toyanswer/.style={draw=black!35,rounded corners=3pt,align=center,
        minimum width=2.4cm,minimum height=1.0cm,font=\normalsize}]
      \node[toyanswer] (a) at (0,0) {$A_1$: $x=4$\\Incorrect};
      \node[toyanswer,draw=ForestGreen,text=ForestGreen] (b) at (4.0,0) {$A_2$: $x=2$\\Correct};
      \node[toyanswer] (c) at (8.0,0) {$A_3$: $x=5$\\Incorrect};
      \draw[->,line width=.8pt,draw=black!60] (a) -- node[above,font=\small] {revise} (b);
      \draw[->,line width=.8pt,draw=black!60] (b) -- node[above,font=\small] {revise} (c);
    \end{tikzpicture}
  \end{center}
  \par\medskip
  {\Large\color{unirot}\textbf{About 38\%}} of correct answers became incorrect under the authors' naive revision approach.
  \par\medskip
  Selecting from the history gives the system a chance to recover an earlier correct answer.
  \note{Optional backup. The three-answer chain is a schematic example, not measured outputs. The reported 38 percent concerns correct answers converted to incorrect ones under naive revision, not 38 percent of all generated answers. This is consistent with a distribution shift: training references were incorrect, while generated inference history can already contain correct answers. Selecting from history can recover a correct answer, but neither verifier selection nor majority voting guarantees correctness. Figure 6 right compares sequential and parallel generation after selection, rather than directly comparing last-answer submission with selection.}
\end{frame}



\begin{frame}{Backup: A Complete Beam-Search Example}
  Suppose each round maintains $N=6$ candidates and each retained prefix gets $M=2$ continuations.
  \par\medskip
  \begin{enumerate}
    \item Generate 6 possible first steps and score their prefixes with the PRM.
    \item Keep the top $N/M=3$ prefixes.
    \item Extend each twice, giving $3\times2=6$ candidates for the next round.
    \item Repeat scoring and expansion until completion or 40 rounds.
    \item Apply weighted best-of-$N$ to the final candidates and return one answer.
  \end{enumerate}
  \par\medskip
  \keyterm{The base model proposes, the PRM scores, search allocates effort, and final selection chooses the output.}
  \note{Optional backup. Recap the complete process after explaining the three components. Pruning happens across all candidate prefixes, so two descendants of one prefix can both survive. N is the candidate pool and M is the number of continuations per retained prefix in this implementation. The general workflow drawing omitted these algorithm-specific settings. Forty is the paper's maximum expansion count. Use this example if a question calls for the detailed expansion and pruning procedure.}
\end{frame}


\begin{frame}{Backup: Difficulty Estimation and Policy Selection}
  Test this idea with \textbf{PaLM 2-S* (Codey)} on \textbf{MATH}: 12,000 training questions and 500 test questions.
  \par\medskip
  \begin{enumerate}
    \item Estimate difficulty from 2,048 independent samples per question.
    \item Sort questions into five equal-frequency bins, from easy to hard.
    \item For each bin and budget, select a setting on one fold and evaluate on the other. Swap folds and average.
  \end{enumerate}
  \par\smallskip
  \begin{tabular}{@{}ll@{}}
    \textbf{Oracle bins} & Rank by actual correctness rate.\\
    \textbf{Predicted bins} & Rank by the verifier's mean final-answer score.
  \end{tabular}
  \par\medskip
  \keyterm{Reported budgets exclude difficulty-estimation cost.}
  \note{Optional backup. With a trained verifier available, return to choosing an allocation by difficulty. Briefly introduce PaLM 2-S* and MATH, with 12,000 training and 500 test questions. Use actual correctness or mean verifier scores to construct two separate five-bin partitions. The latter avoids needing ground-truth answers when scoring difficulty. Choose settings on one fold and evaluate on the other, then swap and average. Difficulty-estimation cost is excluded. We now have a scorer and a selection rule; next explain how its scores determine a final answer.}
\end{frame}


\begin{frame}{Backup: Difficulty Bins and Cross-Validation}
  \begin{enumerate}
    \item Score every question using 2,048 samples. Sort and divide into five quantiles.
    \item Within a bin, split questions into folds A and B.
    \item Use fold A to select settings at each budget. Evaluate those settings on B.
    \item Swap A and B. Average held-out performance.
  \end{enumerate}
  \par\medskip
  Oracle and predicted difficulty repeat this procedure with different scoring rules.
  \papersource{Section 3.2.}
  \note{Optional backup. Oracle uses actual correctness, predicted difficulty uses mean verifier scores. The bin boundaries are quantiles of scores across questions. A score such as 0.7 has no universal fixed bin.}
\end{frame}

\begin{frame}{Backup: Training Rollouts and Lookahead}
  \begin{description}
    \item[PRM training] Sample multiple continuations from a prefix to estimate expected final correctness. This supplies supervision.
    \item[Lookahead search] At inference time, generate up to $k$ extra steps with temperature 0. Use the trained PRM's score at that endpoint to evaluate the current prefix.
  \end{description}
  \par\medskip
  The temporary lookahead steps help choose which current prefixes to retain. Beam search is the $k=0$ case.
  \papersource{Sections 5.1--5.2.}
  \note{Optional backup. These two uses of rollout occur at different stages. Training rollouts create target values, while lookahead consults a trained PRM during search and incurs extra inference cost.}
\end{frame}

\begin{frame}{Backup: Deriving the Inference Multiplier}
  \[
    X+aY=M(X+Y)
  \]
  \[
    a=M+(M-1)\frac{X}{Y}
  \]
  \[
    \frac{X}{Y}=\frac{6PD_{\mathrm{pretrain}}}{2PD_{\mathrm{inference}}}
    =\frac{3}{R}
  \]
  \[
    \boxed{a=M+\frac{3(M-1)}{R}}
  \]
  Approximate FLOPs per token: forward $2P$, backward $4P$, training total $6P$.
  \papersource{Section 7.}
  \note{Optional backup. Backward computation forms gradients for weights and for activations, with each major matrix computation comparable to a forward pass. Total P includes parameters in every layer, so layer count is already accounted for. These approximations omit some operations.}
\end{frame}

\end{document}
